\title{基于XGBoost预测与线性规划模型的微网购电及储能调控策略}
\maketitle


\begin{abstract}
为解决微网和外网在不同信息条件下的购电与储能调控问题，在满足负载需求的同时尽可能降低运行费用，本文构建了基于{\heiti 线性规划模型}的优化框架，结合 {\heiti XGBoost} 等预测方法、预测误差风险修正与滚动优化，建立不同信息条件下的微网购电策略，在满足供需平衡和储能运行要求的同时降低购电费用与紧急购电风险。

{\heiti 针对问题一，}本文综合考虑分时电价、负载、光伏功率及储能在不同时段间转移电能的特性，建立{\heiti 线性规划模型}，联合优化全天各时段的购电与储能充放电。所得方案中，全天计划购电量为 $59,482.70$,kWh，购电费用为 \textbf{$35,126.85$ 元}，弃电量为 $0$，相比{\heiti 分时贪婪策略}，购电费用降低 \textbf{$9.6\%$}，表明综合全天电价与供需信息进行联合优化的方法能够获得更经济的购电与储能调度方案。

{\heiti 针对问题二，}本文采用{\heiti 扩展窗口 XGBoost}预测各时段负载，并以近 $7$ 日历史同一时刻的平均功率预测光伏功率；根据历史净负载预测误差的{\heiti $0.8$ 分位数}修正净负载预测值，并将其代入线性规划模型，制定全天计划购电与储能充放电方案。计划购电量确定后保持不变，实际运行中根据真实负载与光伏功率调度储能，剩余供电缺口通过紧急购电补充。全年 $334$ 天总购电费用为 \textbf{$13,695,403.61$ 元}，紧急购电量为 $87,331.34$,kWh；相比不进行分位修正的方案，全年购电费用降低约 \textbf{$160.9$ 万元}，表明净负载分位修正能够有效降低预测误差带来的高价紧急购电损失。

{\heiti 针对问题三，}题目要求在 $0/6/12/18$ 点更新光伏预报后判断是否需要调整购电计划。本文继承问题二的点预测骨干，把展望窗延长到当天剩余加次日全天，按净负载分位签合同并做{\heiti 日内滚动调整}，偏差按少买 $0.5$ 倍、多买 $1.5$ 倍结算。锁死段分位 $q_L=0.60$ 只在 $1$ 月 $15$--$31$ 日选定。对照不调整方案，本文全年 \textbf{13\,275\,795.85} 元、紧急购电 $58\,355.20$\,kWh，不调整为 $15\,197\,639.94$ 元、紧急 $464\,543.52$\,kWh，费用低约 \textbf{192.2} 万元（\textbf{12.6\%}），紧急电量降 \textbf{87.4\%}，日末贴底由 $138$ 天降至 $1$ 天。分月看 $11$ 个月中 $9$ 个月更省，仅 $4$ 月与 $8$ 月略贵。题目问是否需要 $6/12/18$ 点调整，答案是{\heiti 需要的}。

{\heiti 针对问题四，}针对实时波动电价下未来电价未知的问题，本文利用历史电价构建日前电价预测模型，依次通过{\heiti 历史同刻均值、星期修正和近 $7$ 日残差修正}预测当天各时段电价，并将结果导入问题二、问题三的模型求解。结果表明，日前购电策略全年总费用为 \textbf{$14,413,479.02$ 元}；采用日内滚动优化后，全年总费用降至 \textbf{$14,001,199.26$ 元}，较不进行日内调整的方案降低 \textbf{$12.4\%$}，说明在实时波动电价条件下，结合电价预测与日内更新信息调整购电计划能够有效降低全年购电费用。
{\heiti 针对问题四，}针对实时波动电价下未来电价未知的情况，本文利用历史电价构建日前电价预测模型，以{\heiti 历史同刻均值}刻画日内价格规律，并依次进行星期修正和近 $7$ 日残差修正，得到全天各时段的预测电价。在此基础上，以预测电价作为优化输入，分别结合问题二的日前购电策略与问题三的日内滚动优化策略制定购电及储能计划，并按实际电价计算运行费用。结果表明，两种策略的全年总费用分别为 \textbf{$14,413,479.02$ 元}和\textbf{$14,001,199.26$ 元}；其中日内滚动策略较不进行日内调整的方案降低费用 \textbf{$12.4\%$}，说明在实时电价环境下，根据预测电价制定购电计划并利用日内更新信息持续调整后续决策，能够有效降低全年运行费用。

\keywords{线性规划；XGBoost；净负载分位；滚动优化；因果电价预测}
\end{abstract}

% 官方规范要求正文不设目录；不要加入 \tableofcontents。
\section{问题重述}

\subsection{研究背景}

微网通过协调光伏、储能与外网购电满足本地负载，可利用低价购电、富余光伏充电及高价时段放电降低运行费用\cite{lasseter2002,katiraei2008}。由于负载、光伏及电价随时间变化，购电计划还需应对预测偏差引起的紧急补购与调整费用。本文利用题给数据，在满足供需平衡和储能运行约束的前提下，研究不同信息条件下的购电与充放电策略。

\subsection{问题提出}

\textbf{问题一：}在给定典型日电价、负载与光伏预测的条件下，制定全天购电及储能充放电计划，使日初、日末储电量相同，并最小化购电费用。

\textbf{问题二：}在负载与光伏未知、每日分时电价相同的条件下，每天零点制定当天购电计划；结合实际运行中的紧急补购，给出指定日期及全年购电结果。

\textbf{问题三：}利用每天 $0/6/12/18$ 点发布的光伏预报制定并更新购电策略，综合计划、调整与紧急购电费用，判断引入日内调整是否有必要。

\textbf{问题四：}在波动电价条件下重新求解问题二、三，分别给出日前购电与日内调整策略的计算结果。各问的模型衔接见图~\ref{fig:overall_framework}。

\section{模型假设与符号说明}

\subsection{模型假设}
只列出对模型结论有实际影响且能够辩护的假设，并说明每条假设的作用。

\subsection{符号说明}
\begin{table}[htbp]
  \centering
  \caption{主要符号及其含义}
  \label{tab:symbols}
  \begin{tabularx}{0.9\textwidth}{>{\centering\arraybackslash}p{0.22\textwidth}X>{\centering\arraybackslash}p{0.18\textwidth}}
    \toprule
    符号 & 含义 & 单位 \\
    \midrule
    $x$ & 在此填写变量含义 & 在此填写 \\
    $D$ & 在此填写参数含义 & 在此填写 \\
    \bottomrule
  \end{tabularx}
\end{table}

% =====================================================================
% 数据说明与预处理（精简版）
% ---------------------------------------------------------------------
% 使用方式：替换主文档中的“数据说明与预处理”一节；依赖图片与主文档一致：
%   Q1_price_load_pv.pdf     典型日负载/光伏/电价
%   att2_energy_b.png        全年净负载持续曲线
%   att4_panel_b.png         全年电价日内分布（中位数+四分位带）
%   att2_seasonal_{winter,spring,summer,autumn}.png  四季日内曲线
% 精简说明：原 4 个独立 figure 合并为 1 个 2×2 组合图；四季四张小图用
%   tabular 内嵌排布，避免 subfigure 内 \\ 换行导致的文字重叠；原“全年
%   日电量”（att2_energy_a.png）与“全年电价日历热图”（att4_panel_a.png）
%   信息已折入正文，不再单独成图，以压缩 1–2 页。
% =====================================================================

\section{数据说明与预处理}

\subsection{数据来源与说明}

本题五份附件均为 Excel 格式：附件 1 是典型日分时电价、负载与光伏预测，支撑问题一的确定性调度；附件 2 是全年 365 天负载与光伏实测，用于问题二的滚动预测与回测；附件 3、4 分别是光伏日前预报与全年电价，用于问题三、四；附件 5 是储能参数与结果模板。字段口径见表~\ref{tab:fields}。

\begin{table}[htbp]
\centering
\small
\caption{数据字段说明}
\label{tab:fields}
\begin{tabular}{llll}
\toprule
字段名称 & 字段含义 & 数据类型 & 单位 \\
\midrule
电价 & 分时购电价格（附件1 单日、附件4 全年） & 数值型 & 元/kWh \\
小区负载 & 小区用电功率（附件1、附件2） & 数值型 & kW \\
光伏发电功率 & 光伏出力（附件1 预测、附件2 实测） & 数值型 & kW \\
光伏日前预报 & 按发布时刻给出的光伏预报（附件3） & 数值型 & kW \\
日期 / 时间 & 10 min 间隔的记录时刻 & 时间型 & — \\
\bottomrule
\end{tabular}
\end{table}

\subsection{数据预处理}

逐表核查：附件 1、2、4 的数值字段缺失值、负值、重复行均为 0，附件 2、4 日期连续 365 天；附件 3 数值字段无缺失、无负值，仅“日期”列因每天 4 次预报合并单元格，按前向填充补齐。按逐时段 IQR 检出的离群点（负载 1308、光伏 34、电价 89）逐一核对均落在冬季负载高峰、夏季强光、峰谷尖峰等真实波动上，故不做异常值剔除与标准化。真正需要处理的是时间口径、单位、夜间噪声与面向模型的特征构造：

\begin{enumerate}[label=\arabic*.)]
  \item \textbf{时间对齐}：附件时刻为时段起点（$0{:}00{+}1$ 表示次日 $00{:}00$--$00{:}10$），与结果模板表头一致；优化覆盖当天 $00{:}00$--$24{:}00$ 共 144 段，全年负载/光伏将前一日末段接到当天首段。
  \item \textbf{单位换算}：目标函数与储能递推均用电量（kWh），故按 $E=P\,\Delta t$（$\Delta t=1/6$\,h）把附件功率换算为电量。
  \item \textbf{光伏夜间置零}：全年恒为 0 的 43 个时段强制置 0；其余晨昏过渡段残值最大约 $10$\,kW，相对负载数千 kW 可忽略，保留不处理。
  \item \textbf{附件 3 尺度对齐}：预报按整点给出，与 10 min 实测尺度不一致，按被预报钟点提取整点值与实测对齐。
  \item \textbf{特征构造}：为 XGBoost 构造时刻编号、星期、月份、年内日序，以及前 $1/2/7$ 天同刻值、近 $3/7/30$ 天同刻均值、$7$ 天同刻标准差共 11 个特征，且只用预测日前数据，避免信息泄漏。
  \item \textbf{残差分位保守修正}：点预测偏差会转化为供电缺口，故问题二将净负载残差的 $\alpha=0.8$ 高分位叠加到点预测上（见问题二），使计划购电偏保守、降低紧急购电风险。
\end{enumerate}

\subsection{数据探索性分析}

核心指标描述性统计见表~\ref{tab:desc}。

\begin{table}[htbp]
\centering
\small
\caption{核心指标描述性统计}
\label{tab:desc}
\begin{tabular}{lccccc}
\toprule
指标 & 均值 & 标准差 & 最小值 & 最大值 & 中位数 \\
\midrule
电价（附件1，元/kWh） & 0.77 & 0.31 & 0.37 & 1.40 & 0.81 \\
小区负载（附件1，kW） & 4626.03 & 977.47 & 3309.39 & 5958.97 & 4377.57 \\
光伏预测（附件1，kW） & 2311.78 & 2854.40 & 0 & 7612.32 & 387.61 \\
小区负载（附件2 全年，kW） & 4626.03 & 1411.73 & 1995.72 & 7978.88 & 4270.68 \\
光伏实测（附件2 全年，kW） & 2311.79 & 2948.64 & 0 & 10216.20 & 18.86 \\
\bottomrule
\end{tabular}
\end{table}

由表~\ref{tab:desc} 可见：光伏全年中位数仅 $18.86$\,kW、均值却达 $2311.79$\,kW，二者相差悬殊，说明光伏存在大量夜间零出力时段、白天集中出力；负载全年极差 $1995.72\sim7978.88$\,kW、标准差 $1411.73$\,kW，波动明显大于典型日，须逐日滚动预测而非套用固定模板。图~\ref{fig:eda} 把典型日、全年净负载、四季曲线与全年电价整合为一张图。

\begin{figure}[htbp]
  \centering
  \begin{subfigure}[t]{0.48\textwidth}
    \centering
    \includegraphics[width=\linewidth]{Q1_price_load_pv.pdf}
    \caption{典型日负载、光伏与分时电价}
    \label{fig:eda_typical}
  \end{subfigure}\hfill
  \begin{subfigure}[t]{0.48\textwidth}
    \centering
    \includegraphics[width=\linewidth]{att2_energy_b.png}
    \caption{全年净负载持续曲线}
    \label{fig:eda_netload}
  \end{subfigure}

  \vspace{1.2em}

  \begin{subfigure}[t]{0.48\textwidth}
    \centering
    \includegraphics[width=\linewidth]{att4_panel_b.png}
    \caption{全年电价日内分布}
    \label{fig:eda_price}
  \end{subfigure}\hfill
  \begin{subfigure}[t]{0.48\textwidth}
    \centering
    \begin{tabular}{@{}c@{\hspace{0.6em}}c@{}}
      \includegraphics[width=0.22\textwidth]{att2_seasonal_winter.png} &
      \includegraphics[width=0.22\textwidth]{att2_seasonal_spring.png} \\[0.4em]
      \includegraphics[width=0.22\textwidth]{att2_seasonal_summer.png} &
      \includegraphics[width=0.22\textwidth]{att2_seasonal_autumn.png}
    \end{tabular}
    \caption{四季日内负载与光伏曲线}
    \label{fig:eda_season}
  \end{subfigure}

  \caption{数据探索性分析：(a) 典型日、(b) 全年净负载、(c) 全年电价、(d) 四季曲线}
  \label{fig:eda}
\end{figure}

由图~\ref{fig:eda} 可读出：
\begin{enumerate}[label=\arabic*.)]
  \item 典型日（图 a）电价呈“夜间低谷、早晚双峰、正午次低谷”形态，全天最低 $0.3713$ 元/kWh（05:40）、最高 $1.3952$ 元/kWh（20:40）；光伏仅在 $05{:}00$--$19{:}00$ 出力、正午峰值约 $7600$\,kW，负载 $08{:}00$--$17{:}00$ 维持高位。光伏与电价低点重叠于正午、电价高峰在早晚，为储能“低价充电、高价放电”提供了空间。
  \item 全年净负载持续曲线（图 b）显示净负载为正的时段占 $80.7\%$、范围 $[-6602,6688]$\,kW，正负两端均超过接口功率 $5000$\,kW；日电量上负载相对平稳、光伏夏高冬低、净缺口与光伏大致反向。储能既不能完全消纳光伏过剩，也不能补足最大缺口，模型必须同时设置弃光与紧急购电两个变量。
  \item 全年电价（图 c）日内峰谷明显、范围 $0.0076\sim1.7936$ 元/kWh，其中与附件 1 相关性最高的日期为 2025-04-22（相关系数 $0.992$），说明附件 1 足以代表全年电价形态。
  \item 四季曲线（图 d）显示负载四季差异较小但夜间始终偏高，光伏夏季日均 $2673$\,kW、约为冬季 $1681$\,kW 的 $1.6$ 倍，季节差异显著，光伏预测须引入季节特征。
\end{enumerate}

此外，工作日负载峰值 $6195$\,kW、周末 $5367$\,kW，工作日高约 $15\%$，其星期结构由问题二的负载回归处理；附件 3 的 $0{:}00$ 光伏预报白天平均 MAE 约 $311$\,kW，晴天与多云可差一个数量级，详见问题三。

\subsection{数据处理结果小结}

综上，原始数据完整、无缺失、无负值，异常点均为真实波动予以保留，处理工作集中在时间口径对齐与面向模型的特征构造。处理后得到统一 10 分钟尺度的规范数据集：附件 2、4 各 $365\times144=52560$ 条，附件 1 为 144 条，附件 3 对齐至 10 分钟尺度，并构造出问题二所需的 11 个特征。后续问题一基于附件 1 开展确定性线性规划，问题二基于附件 2 开展滚动预测、日前优化与回测，问题三、四基于附件 3、4 开展不确定性分析与全年电价建模。
% 按实际问题数在 sections/questions.tex 中维护各个 Qx 的引用。
% 按实际赛题的问题数量维护以下引用，不要假设所有比赛只有三问。
% 默认示例（可按实际题目删除或增加）：
\section{问题一：确定性模型}

\subsection{问题分析}
说明问题一的任务类型、输入输出和建模思路。

\subsection{模型建立}
给出决策变量、目标函数、约束条件和必要的推导。每个符号应先在符号表中定义。

\subsection{模型求解与结果}
说明求解方法、参数设置、主要结果及其单位，并用表格或图形展示关键结论。

\subsection{模型检验}
给出基线对照、误差分析或约束可行性检查。

\section{问题二：日前预测购电与日内滚动回测模型}

\subsection{问题分析}

问题二需在每天 $0{:}00$ 仅利用此前历史数据及已知日历信息，制定并锁定当天购电计划。与问题一相比，负载与光伏的预测偏差可能造成供电缺口，须按交易时刻电价的 $5$ 倍紧急补购，因此关键在于权衡日前购电成本与缺电风险。

针对负载的星期周期与光伏的短期连续性，分别采用扩展窗口 XGBoost 和近 $7$ 日同刻均值预测，并以历史净负载残差分位修正预测结果。随后沿用问题一的线性规划框架，加入日末留电要求，联合优化购电与储能运行。最后以真实负载、光伏逐时段执行锁定计划，通过储能预留和紧急补购处理缺口，并将实际日末储电量传递至次日，检验连续运行的费用与可行性。

\subsection{模型建立}

\subsubsection{决策变量}

本问沿用问题一的决策变量及约束条件。由于问题二采用连续逐日运行方式，日前规划不再要求日末储电量恢复至 $6000\,\mathrm{kWh}$，而设置 $E_{144}=S^{*}$，其中 $S^{*}$ 为计划日末储电量，其取值在后文参数标定中确定。

\subsubsection{日前点预测}

每天 $0{:}00$ 制定购电计划时，当天真实负载与光伏功率尚未知，因此仅利用决策日前的历史数据预测当天 $144$ 个时段的负载与光伏功率，分别记为 $\hat P^{\mathrm{load}}_{t}$ 和 $\hat P^{\mathrm{pv}}_{t}$。对第 $D$ 日进行预测时，仅使用日期早于 $D$ 的历史数据，并可使用当天已知的时段、星期、月份等日历信息。由于负载与光伏具有不同的时间规律，本文分别建立预测模型。

\textbf{（1）负载预测：扩展窗口 XGBoost。}

对附件 2 一月负载按星期进行方差分解，星期组间平方和占总平方和的 $0.89$；隔周相关系数为 $0.99$，而隔日相关系数仅为 $0.37$，表明负载具有显著的星期周期。相应地，近 $7$ 日同刻均值、昨日同刻和同星期同刻预测的 MAE 分别为 $781$、$650$ 和 $319\,\mathrm{kW}$。其中，同星期同刻方法虽能利用周周期，但难以进一步反映近期负载变化，因此需要同时引入星期信息与近期历史特征。

基于上述特点，本文采用扩展窗口 XGBoost\cite{chen2016xgboost}。模型以每天 $144$ 个时段的负载功率为预测对象，输入时段、星期、月份等日历特征，以及前 $1/2/7$ 日同刻负载、近 $3/7/30$ 日同刻均值和近 $7$ 日同刻标准差。每天 $0{:}00$ 仅利用此前历史样本重新训练，并预测当天负载。

扩展窗口 XGBoost 的负载 MAE 为 $136.42\,\mathrm{kW}$。进一步比较训练窗口，扩展窗口、滚动 $30$ 日和滚动 $60$ 日的 MAE 分别约为 $136$、$163$ 和 $187\,\mathrm{kW}$，对应全年购电费用约为 $1533$、$1624$ 和 $1696$ 万元。说明保留较长历史能够提供更多有效的周周期样本，因此最终采用扩展窗口 XGBoost。

\textbf{（2）光伏预测：近 $7$ 日同刻均值。}

光伏的星期组间平方和占比仅为 $0.0007$，而相邻日期相关系数达到 $0.96$，说明其星期规律较弱，但短期连续性较强。因此，本文采用近 $W$ 日同刻均值预测第 $D$ 日光伏功率：
\begin{equation}
\hat{P}_{D,t}^{\mathrm{pv}}
=
\frac{1}{W}
\sum_{k=1}^{W}
P_{D-k,t}^{\mathrm{pv}},
\qquad
t=1,\ldots,144
\label{eq:q2_pv_forecast}
\end{equation}

利用 $1$ 月 $15$--$31$ 日数据比较不同窗长，$W=5\sim7$ 时 MAE 最低，约为 $92\sim96\,\mathrm{kW}$；$W=1$ 和 $W=21$ 时均增至约 $119\,\mathrm{kW}$。综合预测精度与近期数据代表性，最终取 $W=7$。

进一步比较发现，跨季节同星期同刻均值的 MAE 为 $493\,\mathrm{kW}$，包含星期特征的树模型为 $159.33\,\mathrm{kW}$，均高于近 $7$ 日同刻均值的 $153.40\,\mathrm{kW}$，进一步说明光伏预测更适合利用近期同刻信息。此外，负载与光伏逐时段 Pearson 相关系数的全年均值仅为 $0.031$，且二者时间规律明显不同，因此分别采用扩展窗口 XGBoost 和近 $7$ 日同刻均值进行预测。

把预测方法与调度环节放在同一张图里对照（图~\ref{fig:q2_model_compare}）。

\begin{figure}[!htbp]
  \centering
  \includegraphics[width=0.86\textwidth]{Q2_fig3_model_compare.png}
  \caption{不同预测与控制策略的全年总购电费对照}
  \label{fig:q2_model_compare}
\end{figure}

由图~\ref{fig:q2_model_compare} 可见：负载与光伏都用 $7$ 日均值时总购电费最高，说明 $7$ 日均值法并不高效；XGB 的滚动窗越长费用越低，说明 XGB 能有效利用历史；本文“负载扩展窗 XGB $+$ 光伏 $7$ 日均值”整体费用最低，叠加后文净分位、日末 $2400$ 与 $\rho$ 后为 $1369.54$ 万元，低于固定分位 $0.8$ 的 $1376.52$ 万元。

\subsubsection{净负载分位与日前线性规划}

点预测只能给出“最可能的净负载”，不能表达误差方向。日前计划的关键不是把净负载估准，而是决定为不确定的净负载预留多少计划电。先看单个时段、把储能耦合暂时拿掉：多买 $1\,\mathrm{kWh}$ 计划电的成本是 $\pi$；若缺口 $1\,\mathrm{kWh}$，须按 $5\pi$ 紧急补购，而这 $1\,\mathrm{kWh}$ 若已买进本只需 $\pi$，故少备的边际损失 $C_u=5\pi-\pi=4\pi$；多备 $1\,\mathrm{kWh}$ 已付 $\pi$，用不上则弃光（残值记 $0$），故多备的边际损失 $C_o=\pi$。临界覆盖率为
\begin{equation}
\alpha^{*}=\frac{C_u}{C_u+C_o}=\frac{4\pi}{4\pi+\pi}=0.8,
\label{eq:q2_critical}
\end{equation}
$\pi$ 约去后与电价表无关。储能能让多购电部分留存，真实 $C_o$ 略小于 $\pi$，故 $0.8$ 是单时段近似下的覆盖率。

因此不在负载、光伏上分别取分位，而在净负载上做保守修正，避免把“负荷上偏、光伏下偏”拆成两次平移。记实测净负载、点预测净负载与残差
\begin{equation}
\begin{aligned}
N_t&=P^{\mathrm{load}}_{t}-P^{\mathrm{pv}}_{t},\\
\hat N_t&=\hat P^{\mathrm{load}}_{t}-\hat P^{\mathrm{pv}}_{t},\\
r^N_{d',t}&=N_{d',t}-\hat N_{d',t},\qquad d'<D.
\end{aligned}
\label{eq:q2_netdef}
\end{equation}
假设给定时段 $t$ 的残差在日与日之间可交换，经验 $0.8$ 分位取决策日前全部历史（expanding）：最近 $7$ 天只有 $7$ 个点，$0.8$ 分位接近第 $6/7$ 大顺序统计量、方差过大；按季切开又使每季开头样本更少。于是取
\begin{equation}
\tilde N_t=\hat N_t+Q_{0.8}(r^N_{\cdot,t}),
\label{eq:q2_net_quantile}
\end{equation}
送入日前 LP 时保持光伏点预测不变（夜间归零），把保守净负载加回负载侧
\begin{equation}
\tilde P^{\mathrm{load}}_{t}=\max\bigl(\hat P^{\mathrm{pv}}_{t}+\tilde N_t,0\bigr).
\label{eq:q2_load_back}
\end{equation}
以修正后的预测覆盖全天净负载，在 $0{:}00$ 求解日前线性规划：
\begin{equation}
\begin{aligned}
\min\quad &\sum_{t=1}^{144}\pi_t G_t \\
\mathrm{s.t.}\quad
&G_t+\tilde P^{\mathrm{pv}}_{t}\Delta t+d_t=\tilde P^{\mathrm{load}}_{t}\Delta t+c_t+w_t,\\
&E_t=E_{t-1}+\eta c_t-d_t/\eta,\\
&E_{\min}\le E_t\le E_{\max},\\
&0\le c_t,d_t\le P_{\max}\Delta t,\qquad G_t,w_t\ge0,\\
&E_{144}=S^{*}=2400\,\mathrm{kWh}.
\end{aligned}
\label{eq:q2_lp}
\end{equation}
其中 $E_0$ 取上一日实际执行末值（$2$ 月 $1$ 日为 $6000\,\mathrm{kWh}$）；日末 $2400\,\mathrm{kWh}$ 只约束计划轨迹，实际执行并不重置到该值，日前 LP 只锁死计划购电 $G_t$，充放电按日内回放重算。

\subsubsection{日内 $\rho$ 保护回放}

“保护”指缺口到来时不把储能一次放到下限，“回放”指计划购电锁死后用当天真值逐时段执行。记净残差
\begin{equation}
r_t=P^{\mathrm{load}}_{t}\Delta t-P^{\mathrm{pv}}_{t}\Delta t-G_t.
\label{eq:q2_resid}
\end{equation}
若 $r_t\le0$，仍按贪心充电、剩余弃光，紧急电不用于充电；若 $r_t>0$，并不把储能放到 $E_{\min}$，而是先按日前计划储电量构造预留线
\begin{equation}
\underline{E}_t=E_{\min}+\rho\bigl(E^{\mathrm{plan}}_{t}-E_{\min}\bigr),
\label{eq:q2_reserve}
\end{equation}
再给出放电上限
\begin{equation}
\min\bigl(r_t,\;\eta(E_{t-1}-\underline{E}_t),\;P_{\max}\Delta t\bigr),
\label{eq:q2_dmax}
\end{equation}
其余记紧急购电 $e_t$，按 $5\pi_t$ 结算。$\rho\in[0,1]$ 是预留比例：$\rho=0$ 即原贪心（放到下限 $1200\,\mathrm{kWh}$），$\rho=1$ 则完全按日前计划储电量放电，越大越不提前把次日要用的库存放完。全天总费用
\begin{equation}
C=\sum_{t}\pi_t G_t+\sum_{t}5\pi_t e_t,
\label{eq:q2_cost}
\end{equation}
回放末值 $E_{144}$ 作为次日实际 $E_0$。

\subsubsection{参数取值与依据}

\begin{enumerate}
  \item \textbf{XGB 超参数。}树深 $6$、学习率 $0.05$、$200$ 棵树、行/列抽样 $0.8$、叶最小权重 $4$、$L_2$ 罚 $1$、种子 $42$；不按 $2$--$12$ 月电费交叉验证，避免用测试期信息选参。
  \item \textbf{光伏窗长 $W=7$。}见日前点预测：$1$ 月 $15$--$31$ 日同刻 MAE 在 $W=5\sim7$ 见底。
  \item \textbf{净分位 $\alpha=0.8$。}由报童临界覆盖率 $4\pi/(4\pi+\pi)=0.8$ 直接给定，残差取决策日前全部历史。
  \item \textbf{日末 $S^{*}=2400$ 与 $\rho=0.625$。}日末留电若回到中位 $6000\,\mathrm{kWh}$，当天谷电大量留给次日；若贴近下限 $1200\,\mathrm{kWh}$，次日早高峰只能靠购电，故 $S^{*}$ 存在折中。$\rho$ 过小会提前把次日要用的电放完，过大则日内缺口几乎全买紧急电。两项参数用 $1$ 月 $15$--$31$ 日做因果标定：以 $1200\,\mathrm{kWh}$ 为步长，在 $S^{*}\in\{2400,3600,4800,6000\}$、$\rho\in\{0.5,0.625\}$、$\alpha\in\{0.8,0.825\}$ 的 $16$ 个组合上，从储电量 $6000\,\mathrm{kWh}$ 起按当天真值反事实回放，并把 $1$ 月 $31$ 日 $24{:}00$ 的储电量折成现金计入总代价，最低的一组是 $S^{*}=2400$、$\rho=0.625$（图~\ref{fig:q2_jan_grid}）。
\end{enumerate}

\begin{figure}[!htbp]
  \centering
  \includegraphics[width=0.9\textwidth]{q2_january_sensitivity_grid.pdf}
  \caption{$1$ 月参数标定网格图}
  \label{fig:q2_jan_grid}
\end{figure}

由图~\ref{fig:q2_jan_grid} 可见：$\alpha=0.825$ 整面比 $\alpha=0.8$ 贵约 $2.7$ 万元，故不取更厚的分位。

每天 $0{:}00$ 重训 XGB，并求解一次含 $144\times5$ 个连续变量的日前线性规划，
$334$ 天均求得最优解且无约束校验错误。全年（$2025$ 年 $2$ 月 $1$ 日至 $12$ 月 $31$ 日）
总购电费为 $\textbf{13\,695\,403.61}$ 元，其中计划购电费 $13\,184\,604.50$ 元、
紧急购电费 $510\,799.11$ 元；紧急购电量 $87\,331.34$\,kWh，
分布在 $137$ 天共 $825$ 个时段；全年计划购电量 $21\,521\,947.29$\,kWh，
弃电量 $2\,472\,688.96$\,kWh。实际储电量始终位于
$[1200,10800]$\,kWh，最低为 $1200$\,kWh、最高为 $10800$\,kWh，
且无同时充放电。$12$ 月 $31$ 日 $24{:}00$ 储电量为 $2678.35$\,kWh。

全年日购电费呈明显的“冬高夏低”特征：$12$ 月日均 $5.46$ 万元最高，
$4$--$5$ 月为 $3.07$--$3.27$ 万元最低；$6$ 月因紧急购电出现全年最高单日费用，
当月紧急购电成本占全年紧急购电成本的 $37.2\%$。

题目指定日期的计划购电与储能调度见表~\ref{tab:q2_purchase}、
表~\ref{tab:q2_storage}。

% =========================================================
% 表6、表7并排
% 通过不同 arraystretch 使上下边界尽量同时对齐
% =========================================================

\begin{samepage}

\begin{table}[H]
  \centering

  % ==================== 左：表6 ====================
  \begin{minipage}[t]{0.49\textwidth}
    \vspace{0pt}
    \centering

    \caption{问题二：指定日期计划购电量及全天购电费}
    \label{tab:q2_purchase}

    \vspace{0.25em}

    \scriptsize
    \setlength{\tabcolsep}{1.2pt}

    \renewcommand{\arraystretch}{1.00}

    \begin{tabular*}{\linewidth}{
      @{\extracolsep{\fill}}
      ccccc
      @{}
    }
      \toprule
      时间段
        & 3.20
        & 6.21
        & 9.23
        & 12.21 \\
      \midrule

      10:00--10:10
        & 0.00
        & 0.00
        & 0.00
        & 0.00 \\

      12:00--12:10
        & 641.27
        & 0.00
        & 554.73
        & 1052.69 \\

      14:00--14:10
        & 0.00
        & 0.00
        & 0.00
        & 0.00 \\

      16:00--16:10
        & 466.08
        & 127.67
        & 475.19
        & 841.17 \\

      18:00--18:10
        & 743.26
        & 437.62
        & 806.06
        & 739.53 \\

      20:00--20:10
        & 0.00
        & 0.00
        & 0.00
        & 0.00 \\

      \midrule

      全天计划购电量（kWh）
        & 70698.38
        & 34922.73
        & 69093.42
        & 99124.96 \\

      全天计划购电费（元）
        & 43184.03
        & 20252.71
        & 42877.32
        & 64123.11 \\

      紧急购电费（元）
        & 0.00
        & 0.00
        & 1933.20
        & 0.00 \\

      全天总购电费（元）
        & 43184.03
        & 20252.71
        & 44810.52
        & 64123.11 \\

      \bottomrule
    \end{tabular*}

  \end{minipage}
  \hfill

  % ==================== 右：表7 ====================
  \begin{minipage}[t]{0.49\textwidth}
    \vspace{0pt}
    \centering

    \caption{问题二：指定日期储能充放电量及起止储电量}
    \label{tab:q2_storage}

    \vspace{0.25em}

    \scriptsize
    \setlength{\tabcolsep}{0.8pt}

    \renewcommand{\arraystretch}{1.10}

    \begin{tabular*}{\linewidth}{
      @{\extracolsep{\fill}}
      ccccc
      @{}
    }
      \toprule

      时间段
        & 3.20
        & 6.21
        & 9.23
        & 12.21 \\

      & \multicolumn{4}{c}{充电量 / 放电量（kWh）} \\

      \midrule

      0:00--4:00
        & 7718.29/49.42
        & 456.07/118.42
        & 7992.53/0.00
        & 7847.31/144.43 \\

      4:00--8:00
        & 969.00/4426.56
        & 542.27/1735.80
        & 722.97/5767.24
        & 1521.31/667.36 \\

      8:00--12:00
        & 5458.12/3265.56
        & 9472.90/0.00
        & 4774.60/2387.94
        & 5135.59/7756.25 \\

      12:00--16:00
        & 4053.33/266.64
        & 0.00/0.00
        & 5581.23/530.20
        & 7180.91/2220.11 \\

      16:00--20:00
        & 111.19/4859.30
        & 89.44/5164.11
        & 123.97/4997.03
        & 0.00/4779.93 \\

      20:00--24:00
        & 1815.79/3681.54
        & 2268.71/3017.50
        & 1948.30/3517.84
        & 1929.15/3622.79 \\

      \midrule

      $0{:}00$ 储电量
        & 3308.16
        & 3436.12
        & 2993.40
        & 3270.22 \\

      $24{:}00$ 储电量
        & 3033.50
        & 3831.65
        & 2911.26
        & 3199.88 \\

      \bottomrule
    \end{tabular*}

  \end{minipage}

\end{table}

\end{samepage}

\vspace{-0.4em}

\noindent

由表~\ref{tab:q2_purchase}、表~\ref{tab:q2_storage} 可见：
计划购电主要集中在光伏回落的 $16{:}00$、$18{:}00$ 时段，
$10{:}00$ 与 $20{:}00$ 则主要由富余光伏和储能放电覆盖；
储能总体呈“凌晨低价充电、早晚高峰放电”的运行规律。
四个指定日日末储电量均高于计划硬约束 $2400$\,kWh。
其中仅 $9$ 月 $23$ 日出现紧急购电，共 $288.88$\,kWh，
主要集中在晚间负荷高峰。

\end{samepage}


% =========================================================
% 指定日调度图
% =========================================================

指定日 2025-03-20 的 $144$ 个 $10\,\mathrm{min}$ 时段调度
如图~\ref{fig:q2-circular}。

\begin{samepage}

\begin{figure}[H]
  \centering
  \includegraphics[width=0.72\textwidth]{q2_circular_dispatch.pdf}
  \caption{指定日 (2025-03-20) 分时电价、净负载、充放电与储电量环形对照}
  \label{fig:q2-circular}
\end{figure}

\vspace{-0.5em}

\noindent

由图~\ref{fig:q2-circular} 可见，储能状态随全天电价和净负载滚动调整，
而非强制恢复至 $6000$\,kWh。该日凌晨低价时段集中充电，
早晚高峰集中放电：$0{:}00$--$4{:}00$ 充电 $7718.29$\,kWh，
$4{:}00$--$8{:}00$ 与 $16{:}00$--$20{:}00$ 分别放电
$4426.56$\,kWh 和 $4859.30$\,kWh。
SOC 始终处于 $[1200,10800]$\,kWh 范围内，
日初为 $3308.16$\,kWh、日末为 $3033.50$\,kWh。
受储能容量上限影响，当日产生弃电 $4.74$\,MWh，
但无紧急购电和缺供。该日全天计划购电 $70\,698.38$\,kWh，
总购电费为 $43\,184.03$ 元。

\end{samepage}


% =========================================================
% 调度环节对照
% =========================================================

\subsection{调度策略对照与效果分析}

在同一套点预测（负载扩展窗 XGBoost、光伏 $7$ 日同刻）
与 $2$ 月 $1$ 日储电量 $6000$\,kWh 下，
对不同分位口径及运行约束进行对照，结果见表~\ref{tab:q2_control}。

\begin{samepage}

\begin{table}[H]
  \centering
  \scriptsize
  \setlength{\tabcolsep}{3.2pt}
  \renewcommand{\arraystretch}{0.95}

  \caption{同一套点预测下不同调度策略的全年费用对照}
  \label{tab:q2_control}

  \begin{tabular}{lrrrr}
    \toprule
    策略
      & 计划购电费（元）
      & 紧急购电费（元）
      & 紧急电量（kWh）
      & 总购电费（元） \\
    \midrule

    无分位
      & 12124833.88
      & 3198303.43
      & 525856.18
      & 15323137.31 \\

    固定分量 $0.8/0.2$
      & 13369773.46
      & 395394.12
      & 66910.43
      & 13765167.58 \\

    净分位 $\alpha=0.8$
      & 13192124.98
      & 522427.80
      & 87329.02
      & 13714552.77 \\

    \midrule

    净分位 $+$ 日末 $2400$
      & 13185526.68
      & 513021.95
      & 82852.85
      & 13698548.63 \\

    净分位 $+$ 日末 $2400$ $+$ $\rho=0.625$
      & 13184604.50
      & 510799.11
      & 87331.34
      & \textbf{13695403.61} \\

    \bottomrule
  \end{tabular}
\end{table}

\vspace{-0.4em}

\noindent

由表~\ref{tab:q2_control} 可见，
在三种分位口径中，净分位策略较固定分量方案全年总费用降低约
$5.06$ 万元，较无分位方案降低约 $160.9$ 万元；
在净分位基础上加入日末 $2400$\,kWh 约束后，总费用进一步降低约
$1.6$ 万元，再加入 $\rho=0.625$ 后继续降低 $3145$ 元。
因此，最终采用“净分位 $+$ 日末 $2400$\,kWh $+$ $\rho=0.625$”
作为正式调度方案。回测区间逐日费用及分月分布见附录图~\ref{fig:app_q2_daily}、图~\ref{fig:app_q2_month}。

\end{samepage}

% ======================================================================
% 六、问题三：日内滚动购电与调整策略的建模与求解（7 面版）
% 对应仓库：Sen-illion/math-modeling @ task/integrate-q123
% 冻结口径：LA = 48h 展望 + 三段分源残差分位修正（q_L=0.60 一月定参）
%           正式窗 2025-02-01--12-31，2/1 0:00 SOC=6000 kWh
% 需用宏包：amsmath, amssymb, booktabs, graphicx, caption, enumitem
% 图文件：fig_q3_horizon.png, fig_q3_dispatch_0621.png
% ======================================================================
\section{问题三：日内滚动购电与调整策略的建模与求解}

\subsection{问题分析与建模思路}

问题三在问题二“日前一次决策”的基础上，把一天拆成四个决策时点：$0{:}00$ 制定全天计划购电 $G^{\mathrm{plan}}$，随后在 $6{:}00$、$12{:}00$、$18{:}00$ 获得更新的光伏预报后，对尚未执行的合同做增购或减购调整，最终执行 $G^{\mathrm{adj}}$。调整并非免费：比计划少买按 $0.5\pi$ 结算、多买按 $1.5\pi$ 结算，仍出现缺口则按 $5\pi$ 紧急购电。这一问的难点不在“再解一次优化”，而在三件事环环相扣，我们先把它们拆开。

\textbf{（1）信息层——预报如何进入决策。} 附件 3 只给出四个发布时点之后 $T{+}1,\dots,T{+}24$ 的整点光伏预报，既不是 $10$ 分钟粒度，也没有“发布当刻”的起点值。若按整点硬拼，夏季 $6{:}00/18{:}00$ 起点会被系统压成零、低估全天光伏；若误用未发布的次日预报又造成信息泄漏。因此第一件事是\textbf{无泄漏地把整点预报补成 $10$ 分钟序列}，否则后面所有优化都建立在错误输入上。

\textbf{（2）结算层——不对称惩罚决定“分位怎么取”。} 少买 $0.5\pi$、多买 $1.5\pi$、缺口 $5\pi$ 三种价格不对称，直接改变“买多少保险”的答案：两次调整之间\textbf{锁死}的时段下一时刻前无法再改、少买代价高，应按高覆盖率买足；下一时刻\textbf{还能改}的时段，偏差可被下一轮修正，取中位即可。这正是后文“三段分位”的由来。

\textbf{（3）物理层——真正的失效在晚间电量预算。} 储能可用区间只有 $[1200,10800]$ kWh、接口 $5000$ kW。诊断（见 6.8 节）显示，旧方案紧急购电 $88.8\%$ 来自\textbf{晚间储能电量耗尽}而非功率触顶，$20{:}00$ 单小时占 $47.4\%$，日末有 $106$ 天贴近下限。模型必须在优化阶段就为晚间留够电量，而不是事后用 $5\pi$ 补缺口。

综合三层约束，本问被定位为一个\textbf{带追索权（recourse）的多阶段随机线性规划}：第一阶段在预报不确定时先承诺合同，后续阶段看到新预报后再做有代价的修正。目标与约束关于购电量、充放电量都是线性的，不确定性只通过“预报残差”进入，故用“分位点预报”把随机性近似掉，得到可逐时段求解的\textbf{滚动时域线性规划}；不必上随机规划或元启发式，因为随机性来源单一、惩罚结构明确，分位修正已能把风险压到计划阶段。

\subsection{模型假设}

本问沿用前两问公共假设：充放电单程效率各为 $\eta=0.9$（交流侧口径）、储能运行窗 $[1200,10800]$ kWh、接口最大功率 $5000$ kW、不向外网倒送电、光伏过剩允许弃光；一天仍为 $144$ 个 $10$ 分钟时段，$\Delta t=1/6$ h。相对问题二另有三条约定：光伏决策只用附件 3 预报并线性插值到 $10$ 分钟；负荷 $0{:}00$ 沿用因果 Expanding XGBoost、$6/12/18$ 时点用已发生均值比例校正；不强制 $E_0=E_{24}$，储电量跨日传递，正式窗 $2$ 月 $1$ 日 $0{:}00$ 置 $6000$ kWh、$1$ 月仅用于预报与定参。新增符号 $G^{\mathrm{plan}},G^{\mathrm{adj}},\Delta^{\pm},G^{\mathrm{em}},q_L,q_O,q_E$ 在正文就近定义。

\subsection{结算规则与目标函数}

设一天 $144$ 段，$t=1,\dots,144$，交易电价 $\pi_t$。定义上调量与下调量
\begin{equation}
  \Delta^{+}_t=\max\!\big(G^{\mathrm{adj}}_t-G^{\mathrm{plan}}_t,\,0\big),\qquad
  \Delta^{-}_t=\max\!\big(G^{\mathrm{plan}}_t-G^{\mathrm{adj}}_t,\,0\big).
  \label{eq:q3_dpm}
\end{equation}
一天总费用为计划购电费、调整偏差费与紧急购电费之和：
\begin{equation}
  C = \sum_{t=1}^{144}\pi_t G^{\mathrm{plan}}_t
      -\dfrac{1}{2}\sum_{t=1}^{144}\pi_t \Delta^{-}_t
      +\dfrac{3}{2}\sum_{t=1}^{144}\pi_t \Delta^{+}_t
      +5\sum_{t=1}^{144}\pi_t G^{\mathrm{em}}_t.
  \label{eq:q3_cost}
\end{equation}
三项对应题面三档价格：已承诺合同按 $\pi_t$ 计价；少买“退回一半电费”（$-\tfrac12\pi_t\Delta^{-}_t$），多买按 $1.5$ 倍计价（$+\tfrac32\pi_t\Delta^{+}_t$）；真实缺口按 $5\pi_t$ 紧急购电。式~\eqref{eq:q3_cost} 是本问所有策略的统一评价指标。

\subsection{物理约束与储能动态}

以最终执行量 $G^{\mathrm{adj}}$ 替代计划量并计入紧急购电，时段电量平衡为
\begin{equation}
  G^{\mathrm{adj}}_t + P^{\mathrm{pv}}_t\Delta t + d_t + G^{\mathrm{em}}_t
  = P^{\mathrm{load}}_t\Delta t + c_t + w_t,
  \label{eq:q3_balance}
\end{equation}
储能递推与运行约束沿用前两问：
\begin{equation}
  E_t = E_{t-1}+\eta c_t - \dfrac{d_t}{\eta},\qquad
  1200\le E_t\le 10800,\qquad 0\le c_t,\,d_t\le P_{\max}\Delta t,\qquad w_t\ge 0,\qquad G^{\mathrm{em}}_t\ge 0.
  \label{eq:q3_phys}
\end{equation}
其中 $c_t,d_t$ 为交流侧充放电量，$w_t$ 为弃光，$E_t$ 为时段末储电量。这组约束\textbf{既用于优化、也用于回放}，保证“计划”与“结算”口径一致：优化时代入分位修正后的预报，回放时代入实测值。

\subsection{滚动时域线性规划与三段分位修正}

\subsubsection{$0{:}00$ 计划阶段}

每天 $0{:}00$ 用分位修正后的净负荷预报求解当天 $144$ 段线性规划，得到计划合同 $G^{\mathrm{plan}}$：
\begin{equation}
  \min_{G,c,d,w,E,G^{\mathrm{em}}}
  \sum_{j=1}^{144}\pi_j G_j + 5\sum_{j=1}^{144}\pi_j G^{\mathrm{em}}_j
  + \varepsilon\sum_{j=1}^{144}\big(c_j+d_j\big),
  \label{eq:q3_plan}
\end{equation}
约束为式~\eqref{eq:q3_balance}--\eqref{eq:q3_phys} 在修正预报下的形式。$G^{\mathrm{em}}_j$ 是 LP 内按 $5\pi$ 计价的\textbf{虚拟}紧急购电、只为保证可行性，$\varepsilon$ 是打破“同时充放”的极小正则项。

\subsubsection{$\tau>0$ 调整阶段与 48 小时展望}

在时点 $\tau\in\{6,12,18\}$，把当天已承诺合同 $g_j$ 放进“当天剩余 $+$ 次日全天”的 $48$ h 展望窗内重新决策。对当天剩余 $n_{\mathrm{today}}$ 段在 $g_j$ 上做有代价增减，对次日展望段按普通电价重新决策：
\begin{equation}
  \min_{\substack{G,\delta^{\pm},c,d,\\w,E,G^{\mathrm{em}}}}
  \sum_{j=1}^{n_{\mathrm{today}}}\Big(-\tfrac12\pi_j\delta^{-}_j+\tfrac32\pi_j\delta^{+}_j\Big)
  +\sum_{j=n_{\mathrm{today}}+1}^{n}\pi_j G_j
  +5\sum_{j=1}^{n}\pi_j G^{\mathrm{em}}_j,
  \label{eq:q3_adj}
\end{equation}
\begin{equation}
  \mathrm{s.t.}\quad G_j = g_j + \delta^{+}_j - \delta^{-}_j,\ \ \delta^{\pm}_j\ge 0\ (j\le n_{\mathrm{today}});
  \quad \text{式~\eqref{eq:q3_balance}--\eqref{eq:q3_phys}};\ \ E_0 = \mathrm{SOC}_{\tau}.
  \label{eq:q3_adj_con}
\end{equation}
每次解出后，当天剩余合同更新为新最优值，$6/12/18$ 重复即构成滚动决策（图~\ref{fig:q3_horizon}）。

\begin{figure}[!htbp]
  \centering
  \includegraphics[width=0.8\linewidth]{fig_q3_rolling.png}
  \caption{滚动时域决策示意（以 $6{:}00$ 为例）：紧邻 6 小时为锁死段（$q_L=0.60$），当天剩余为可改段（$q_O=0.50$），次日为展望段（$q_E=0.50$）}
  \label{fig:q3_horizon}
\end{figure}

\textbf{为什么拉到 48 小时？} 若只在“当天剩余”内优化，$0{:}00$ 与 $18{:}00$ 的解都会在日末把储电量耗尽、次日清晨被动；把窗延长到“当天剩余 $+$ 次日全天”后，LP 会主动为次日留储电余量，全年日末贴底从 $106$ 天压到 $1$ 天以内。次日负荷与光伏尾段用“目标日前一周同星期几”的已收盘数据，不用尚未发布的次日 $0{:}00$ 预报，避免泄漏。

\subsubsection{三段分源残差分位修正}

点预报只能给出“最可能”的净负荷，不能表达误差方向。记净负荷 $L_t=P^{\mathrm{load}}_t-P^{\mathrm{pv}}_t$，残差 $r_t=L_t-\hat L_t$ 按“同一发布时点、同一被预报钟点”分桶，得经验分布 $\hat F$。对负荷与光伏\textbf{分源}施加覆盖水平：
\begin{align}
  \tilde P^{\mathrm{load}}_t &= \max\!\Big(\hat P^{\mathrm{load}}_t + \hat F^{-1}_{\mathrm{load}}(q),\ 0\Big), \label{eq:q3_pload}\\[2pt]
  \tilde P^{\mathrm{pv}}_t  &= \max\!\Big(\hat P^{\mathrm{pv}}_t  + \hat F^{-1}_{\mathrm{pv}}(1-q),\ 0\Big). \label{eq:q3_ppv}
\end{align}
负荷取 $q$ 分位（向上抬）、光伏取 $1-q$ 分位（向下压），即“宁多备、勿少备”。分位水平按段属性取三档：锁死段 $q_L$、可改段 $q_O=0.50$、次日段 $q_E=0.50$。

\subsubsection{分位水平的报童论证}

锁死段 $q_L$ 由报童临界分位给出\textbf{搜索上界}：少买 $1$ kWh 的缺货净成本 $c_u=5\pi-\pi=4\pi$，多买 $1$ kWh 的过量成本 $c_o=\pi$，故
\begin{equation}
  q^{*}=\dfrac{c_u}{c_u+c_o}=\dfrac{4\pi}{4\pi+\pi}=0.80.
  \label{eq:q3_crit}
\end{equation}
但它只能作上界：锁死段缺额往往先由电池放电兜底、不必然触发 $5\pi$，真实 $c_u<4\pi$、临界分位应下移，故 $q_L$ 在 $\{0.50,\dots,0.75\}$ 内做一次因果回测定下（见 6.6 节）。可改段 $q_O=0.50$ 由结算对称性直接给出：多承诺净代价 $\pi-0.5\pi=0.5\pi$、少承诺净代价 $1.5\pi-\pi=0.5\pi$，两侧相等，中位数最优。

\subsection{求解与参数选择}

\textbf{（1）输入准备。} 光伏预报以发布时点对应横坐标 $0$、$T{+}k$ 对应 $k$，对 $24$ 个整点线性插值到 $10$ 分钟；插值起点取“发布前最后一个已完成时段的实测光伏”（$0{:}00$ 发布时为 $0$），避免夏季 $6/18$ 点起点硬设为零的系统低估。附件 3 的 $0{:}00$ 发布全预见期平均 MAE 约 $341$ kW、误差集中在日照时段，这正是分位余量“省不掉”的证据。

\textbf{（2）锁死段分位 $q_L$ 的一月定参。} 若在全年 $2$--$12$ 月样本上选 $q_L$，谷点在 $0.65$，但那是“看着答案挑参数”。为避免评分年选参，$q_L$ 只在 $1$ 月 $15$--$31$ 日选定：令 $\nu=\min_t \pi_t/\eta=0.4126$ 元/kWh（最低电价买入并储存 $1$ 个可用 kWh 的成本），取评分
\begin{equation}
  J(q_L)=C_{\mathrm{Jan}}(q_L)+\nu\big(6000-E^{\mathrm{Jan}}_{\mathrm{end}}(q_L)\big),
  \label{eq:q3_score}
\end{equation}
即“一月现金费用 $+$ 把月末储电相对 $6000$ kWh 的差额按 $\nu$ 折成现金”。表~\ref{tab:q3_jan} 中 $q_L=0.60$ 评分最低，故冻结 $q_L=0.60$，$2$--$12$ 月只确认、不参与挑选。

\begin{table}[!htbp]
  \centering\footnotesize
  \caption{一月（1 月 15--31 日）$q_L$ 定参网格：$q_L=0.60$ 评分最低}
  \label{tab:q3_jan}
  \begin{tabular}{@{}l r r r r@{}}
    \toprule
    $q_L$ & 一月现金费用/元 & 一月末 SOC/kWh & 评分 $J$/元 & 一月紧急/kWh \\
    \midrule
    \textbf{0.60} & 870,900.68 & 8377.37 & \textbf{869,919.88} & 3,433.65 \\
    0.63 & 871,055.16 & 8413.64 & 870,059.40 & 2,760.25 \\
    0.65 & 871,993.29 & 8440.72 & 870,986.36 & 2,310.19 \\
    0.55 & 872,521.73 & 8273.82 & 871,583.66 & 4,701.13 \\
    0.70 & 876,219.18 & 8555.60 & 875,164.86 & 1,094.51 \\
    \bottomrule
  \end{tabular}
\end{table}

\textbf{（3）两条关键杠杆的消融。} 以 $24$ h 展望 $+\beta$ 基线 $13{,}458{,}730.30$ 元为锚，全年 $334$ 天因果回测（表~\ref{tab:q3_ablation}）表明：48 h 展望单独省 $3.75$ 万元、日末贴底从 $106$ 天降到 $0$；分位修正单独省 $16.46$ 万元；二者叠加约省 $19.48$ 万元。分位修的是\textbf{晚间电量预算}缺口，不是用更多紧急电换账面费用——加入后费用与紧急电量同向变好。

\begin{table}[!htbp]
  \centering\footnotesize
  \caption{全年消融：48 h 展望与分位修正可叠加（以 24 h$+\beta$ 为锚）}
  \label{tab:q3_ablation}
  \begin{tabular}{@{}l r r r r@{}}
    \toprule
    变体 & 总费用/元 & 较锚/元 & 紧急/kWh & 贴底天数 \\
    \midrule
    24 h$+\beta$（锚） & 13,458,730.30 & 0 & 82,228.36 & 106 \\
    48 h$+\beta$ & 13,421,219.84 & $-37{,}510$ & 70,433.46 & 0 \\
    24 h$+q{=}0.65$ & 13,294,093.95 & $-164{,}636$ & 44,783.43 & 37 \\
    48 h$+q{=}0.65$ & 13,263,910.99 & $-194{,}819$ & 39,678.65 & 1 \\
    \bottomrule
  \end{tabular}

  \vspace{2pt}
  \footnotesize 注：本表仅用于分离两条杠杆的贡献，初值为 1 月滚动口径；正式数字以表~\ref{tab:q3_cmp} 的 2 月 1 日 SOC$=6000$ 口径为准。
\end{table}

\subsection{求解结果}

全年（2025-02-01 至 2025-12-31，共 $334$ 天）逐日回放，三种候选结果见表~\ref{tab:q3_cmp}。

\begin{table}[!htbp]
  \centering\footnotesize
  \caption{全年总费用：LA 最优，较不调整的 N0 省 192.2 万元}
  \label{tab:q3_cmp}
  \begin{tabular}{@{}l r r r r r@{}}
    \toprule
    策略 & 全年总费用/元 & 计划购电费/元 & 紧急购电费/元 & 紧急电量/kWh & 调整天数 \\
    \midrule
    \textbf{LA（赢家）} & \textbf{13,275,795.85} & 12,337,110.40 & 342,008.68 & 58,355.20 & 334 \\
    M0L（选择性调整） & 14,076,756.59 & 12,344,922.94 & 1,301,484.79 & 204,224.31 & 182 \\
    N0（不调整） & 15,197,639.94 & 12,342,385.85 & 2,855,254.09 & 464,543.52 & 0 \\
    \bottomrule
  \end{tabular}
\end{table}

表~\ref{tab:q3_cmp} 读出三条结论。其一，\textbf{日内调整必要}：不调整的 N0 全年 $15{,}197{,}639.94$ 元，比 LA 高 $192.2$ 万元（$+14.5\%$），紧急电量约为 LA 的 8 倍，只用 $0{:}00$ 预报扛不住 24 小时内的预报漂移。其二，\textbf{每次都调优于选择性调}：M0L 只在 $334$ 天里调整 $182$ 天，比 N0 省 $112$ 万元、仍比 LA 贵 $80$ 万元，选择性门漏掉“幅度不大但方向正确”的小修正。其三，LA 的紧急电量只有 N0 的 $12.6\%$，把风险压在计划阶段而非事后用 $5\pi$ 买单。

题目指定日期的计划购电与紧急购电见表~\ref{tab:q3_day}；储能指定日期分时段充放电量（题面表 2 格式）及完整 144 段结果见 result3.xlsx。

\begin{table}[!htbp]
  \centering\footnotesize
  \caption{指定日期计划购电量与紧急购电（题面表 1 与表 3 格式）}
  \label{tab:q3_day}
  \begin{tabular}{@{}l r r r r@{}}
    \toprule
    时间段 & 2025.3.20 & 2025.6.21 & 2025.9.23 & 2025.12.21 \\
    \midrule
    \multicolumn{5}{l}{\textbf{计划购电量 / kWh}} \\
    10:00--10:10 & 0.00 & 0.00 & 0.00 & 0.00 \\
    12:00--12:10 & 523.80 & 0.00 & 357.38 & 1006.90 \\
    14:00--14:10 & 0.00 & 0.00 & 0.00 & 0.00 \\
    16:00--16:10 & 620.65 & 55.53 & 505.11 & 836.82 \\
    18:00--18:10 & 715.88 & 411.32 & 789.37 & 719.13 \\
    20:00--20:10 & 0.00 & 0.00 & 0.00 & 0.00 \\
    全天购电/kWh & 61909.72 & 40973.80 & 66868.91 & 95888.88 \\
    全天购电费/元 & 40736.96 & 22107.38 & 45502.61 & 64039.30 \\
    \midrule
    \multicolumn{5}{l}{\textbf{紧急购电（时间段: 购电量/kWh）}} \\
    \multicolumn{5}{l}{2025.3.20：10:00--10:10 31.93；20:40--20:50 131.05} \\
    \multicolumn{5}{l}{2025.6.21：无紧急购电} \\
    \multicolumn{5}{l}{2025.9.23：8:50--10:00 446.61；20:30--20:50 143.22；21:00--21:10 60.15} \\
    \multicolumn{5}{l}{2025.12.21：8:00--8:20 23.68；10:40--11:00 254.38；20:40--20:50 43.38} \\
    \bottomrule
  \end{tabular}
\end{table}

夏季代表日（2025-06-21）滚动执行过程见图~\ref{fig:q3_dispatch}：净负荷晚高峰 $4004.9$ kW@$19.2$ h 由储能放电削峰，日内调整较计划少买 $127$ kWh。

\begin{figure}[!htbp]
  \centering
  \includegraphics[width=0.7\linewidth]{fig_q3_dispatch_0621.png}
  \caption{夏季代表日（2025-06-21）滚动购电策略：储能晚高峰放电削峰，调整购电较计划少买 127 kWh}
  \label{fig:q3_dispatch}
\end{figure}

\subsection{模型检验与稳健性}

\textbf{（1）独立复算。} 导出 result3.xlsx 后由独立脚本 \texttt{check\_result3.py} 从单元格与附件 1 电价重算逐日合同费、偏差结算、SOC 平衡与跨日连续，不复用滚动代码：正式结果恰好 $334$ 天，储电量全程 $[1200,10800]$、功率不超 $5000$ kW，$2$ 月 $1$ 日 $0{:}00$ 为 $6000$ kWh、相邻日首尾相接，全程无泄漏错误。

\textbf{（2）与先知下界的距离。} 阶段 1（2 月 1--14 日）先知（用实测净负荷解单日 LP 的完美信息下界）总费用 $585{,}705.84$ 元，LA 为 $636{,}942.33$ 元，仅高 $8.7\%$；同期 N0 为 $653{,}094.43$ 元。LA 已逼近先知下界，剩余差距主要来自预报残差而非模型结构。

\textbf{（3）参数稳健性。} $q_L$ 在 V 形谷底附近不敏感：$0.60$ 与 $0.65$ 全年只差约 $0.08\%$，不存在“刀尖上的最优”。一月定参 $q_L=0.60$ 在 $2$--$12$ 月确认重跑中通过门禁、与因果重跑逐位一致。

\textbf{（4）失效模式诊断。} 对旧 $24$ h$+\beta$ 方案逐段回放表明，紧急购电 $88.8\%$ 来自\textbf{晚间储能电量耗尽}而非功率触顶，$20{:}00$ 单小时占 $47.4\%$。这解释了“加大光伏余量”与“回放留 SOC”为何无效——都没碰到“晚间电量预算”这个真约束；新方案把修正做进 LP 本身，日末贴底降到 $1$ 天。

\subsection{本问结论}

本问需要引入 $6{:}00$、$12{:}00$、$18{:}00$ 时点预报调整购电：不调整的 N0 全年 $15{,}197{,}639.94$ 元，比最优策略 LA 高 $192.2$ 万元（$+14.5\%$），说明日内调整必要；每次都调优于选择性调（M0L 比 LA 贵 $80$ 万元）。最优策略为 LA——$48$ h 展望加三段分源残差分位修正（$q_L=0.60$，一月定参），全年总费用 $13{,}275{,}795.85$ 元，紧急电量较 N0 下降 $87.4\%$。完整策略见 result3.xlsx。

\section{问题四：实时波动电价下的购电策略}

\subsection{问题分析}

问题四要求在外网电价实时波动的条件下，重新制定问题二与问题三中的购电与储能调度策略。与前几问不同，当天真实电价在制定购电计划时尚未完全获得，而实际费用又需按照实时电价结算，因此本问首先需要解决未来电价未知条件下的价格预测问题，并在此基础上重新进行购电优化。

针对图~\ref{fig:att4} 所示的日内峰谷结构与逐日价格波动，本文根据决策日前已经获得的历史电价，采用历史同刻均值、星期修正和近期修正相结合的方法预测当天各时段电价，以保证预测过程不使用未来信息。随后分别沿用问题二和问题三的模型，根据预测电价以及负载、光伏预测结果制定购电与储能调度计划。最后，利用附件~4 的实际电价计算全年实际运行费用，通过总费用、紧急购电量及不同策略间的费用差异评价购电策略。

\subsection{日前电价预测模型建立}

记第 $d$ 日时段 $t$ 的真实电价与预测电价为 $p_{d,t}$、$\hat p_{d,t}$。利用决策日前的历史数据，将预测分为同刻基础价格、星期修正与近期修正三部分：
\begin{equation}
  \hat p_{d,t}=b_{d,t}+\Delta_{\mathrm{week},d,t}+\Delta_{\mathrm{recent},d,t},
  \qquad t=1,\ldots,144.
  \label{eq:q4_price_forecast}
\end{equation}
设 $n_d$ 为此前历史日数，$w(d)$ 为星期编号，$\mathcal W_d=\{j<d\mid w(j)=w(d)\}$，$m_d=\min(7,n_d)$，则
\begin{align}
  b_{d,t}&=\frac{1}{n_d}\sum_{j<d}p_{j,t}, \label{eq:q4_price_base}\\
  \Delta_{\mathrm{week},d,t}&=\frac{1}{|\mathcal W_d|}\sum_{j\in\mathcal W_d}r_{j,t}, \label{eq:q4_week_correction}\\
  \Delta_{\mathrm{recent},d,t}&=\frac{1}{m_d}\sum_{k=1}^{m_d}e_{d-k,t}. \label{eq:q4_recent_correction}
\end{align}
其中，$r_{j,t}=p_{j,t}-b_{j,t}$、$e_{j,t}=p_{j,t}-b_{j,t}-\Delta_{\mathrm{week},j,t}$ 为各历史日当时形成的两层残差；无相应历史样本时，修正项取零。

$1$ 月 $1$ 日以附件~1 电价初始化 $b_{1,t}$，两项修正均取零；此后每日先预测，待真实电价获得后再更新历史价格与残差。$1$ 月仅用于积累历史，不计入正式费用；$2$ 月起继续按相同规则逐日更新。

\subsection{模型求解}

本问继承问题二、三的预测方法、调度参数及物理约束，将式~\eqref{eq:q2_lp}、\eqref{eq:q3_lp0}、\eqref{eq:q3_lp_adj} 中的决策电价替换为 $\hat p_{d,t}$。正式回放仍从 $2$ 月 $1$ 日 $6000$ kWh 初始储电量出发，$1$ 月储电量不滚入。

Q4-2 在每日 $0{:}00$ 制定合同并保持日内不变，按问题二的规则回放。Q4-3 在 $6/12/18$ 时点根据更新的负载、光伏信息调整合同，但电价始终采用当天 $0{:}00$ 的预测；次日展望段按相同日内时刻延伸该预测电价。

两模型均调用 CBC 逐日求解，各次优化均获得最优解。实际费用分别按式~\eqref{eq:q2_cost}、\eqref{eq:q3_cost} 结算，其中电价取附件~4 的真实值 $\pi_{d,t}=p_{d,t}$，从而区分决策所用预测价与实际结算价。

\subsection{求解结果}

通过上述预测与调度过程，分别得到实时波动电价下问题二的日前购电方案和问题三的日内滚动购电方案。下面分别对 Q4-2 和 Q4-3 的求解结果进行分析。

% 避免三级标题孤悬于页底。
\Needspace{3\baselineskip}
\subsubsection{Q4-2 日前购电结果}

Q4-2 日前购电方案全年总费用为 $14\,413\,479.02$ 元。为进一步分析该方案的全年运行特征，图~\ref{fig:q4_dispatch} 给出了日初储电量、日购电量、弃电量、日成本以及紧急购电情况。

\begin{figure}[!htbp]
  \centering
  \includegraphics[width=0.78\textwidth]{q4_official_dispatch_details.png}
  \caption{Q4-2 日前购电方案日级运行情况}
  \label{fig:q4_dispatch}
\end{figure}

由图~\ref{fig:q4_dispatch} 可见，Q4-2 日前购电方案的运行结果具有明显的季节变化。图(a)中日初实际储电量始终保持在允许范围内，并围绕计划日末储电量目标波动，表明储能能够保持跨日连续运行。图(b)中日购电量与弃电量随季节发生明显变化，反映了全年负载与光伏发电水平变化对外网购电需求的影响。

由图(c)可见，日成本主要由计划购电费用构成，紧急购电费用在大部分日期较低，仅在少数日期明显增加。图(d)进一步表明，发生紧急购电的日期中，应急购电时段数增加时，应急购电成本总体也随之上升。因此，Q4-2 日前购电方案能够满足全年大部分时段的供电需求，但由于全天计划在 $0{:}00$ 提前制定，当负载或光伏预测出现较大偏差时，少数日期仍需要通过紧急购电弥补供电缺口。


% 避免三级标题孤悬于页底。
\Needspace{3\baselineskip}
\subsubsection{Q4-3 日内滚动购电结果}

为分析日内新增信息对购电策略的改善作用，设置三种方案进行比较。其中，N0 为日前固定方案，即每日 $0{:}00$ 制定购电计划后不再进行日内调整；M0L 为选择性日内调整方案，仅在满足调整条件时修改后续尚未执行的购电计划；LA 为本文采用的日内滚动调整方案，在 $6{:}00$、$12{:}00$ 和 $18{:}00$ 根据新获得的负载和光伏信息重新调整后续购电计划。三种方案的全年费用与紧急购电情况见表~\ref{tab:q4_all}。

\begin{table}[!htbp]
  \centering
  \caption{波动电价下问题三的全年结果比较}
  \label{tab:q4_all}
  \begingroup
  \small
  \setlength{\tabcolsep}{5pt}
  \begin{tabular}{@{}lrrr@{}}
    \toprule
    方案 & 总费用（元） & 紧急购电量（kWh） & 较 N0 节省（元） \\
    \midrule
    N0（不调整）
      & $15\,983\,608.51$ & $466\,576.26$ & --- \\
    M0L（选择性调整）
      & $14\,826\,830.64$ & $218\,215.59$ & $1\,156\,777.87$ \\
    \textbf{LA（本文）}
      & \textbf{$14\,001\,199.26$} & \textbf{$67\,153.26$} & \textbf{$1\,982\,409.26$} \\
    \bottomrule
  \end{tabular}

  \par\smallskip
  \begin{minipage}{0.94\textwidth}
    \footnotesize\textit{注：}LA 的合同购电费用为 $12\,979\,047.00$ 元，
    紧急购电费用为 $370\,783.78$ 元，其余费用来自少买与多买偏差结算。
    各方案均采用附件~4 的真实电价按照相同规则计算实际费用。
  \end{minipage}
  \endgroup
\end{table}

由表~\ref{tab:q4_all} 可见，LA 的全年总费用为 $14\,001\,199.26$ 元，较不进行日内调整的 N0 减少 $1\,982\,409.26$ 元，降幅约为 $12.40\%$；紧急购电量由 $466\,576.26$ kWh 降至 $67\,153.26$ kWh，降幅约为 $85.6\%$。与 M0L 相比，LA 仍进一步节省 $825\,631.39$ 元，降幅约为 $5.57\%$。

结果表明，仅依靠 $0{:}00$ 制定的日前计划难以完全应对当天负载和光伏预测误差。LA 在 $6{:}00$、$12{:}00$ 和 $18{:}00$ 利用新获得的信息调整尚未执行的购电计划后，能够及时修正原计划与实际供需之间的偏差，从而显著减少高价紧急购电并降低全年总费用。

\begin{figure}[!htbp]
  \centering
  \begin{subfigure}[t]{0.96\textwidth}
    \vspace{0pt}\centering
    \includegraphics[width=\linewidth]{q4_cost_structure.pdf}
    \caption{全年成本结构}
    \label{fig:q4_cost}
  \end{subfigure}\\[-6pt]
  \begin{subfigure}[t]{0.96\textwidth}
    \vspace{0pt}\centering
    \includegraphics[width=\linewidth]{q4_daily_cost_distribution.pdf}
    \caption{日总成本分布}
    \label{fig:q4_daily}
  \end{subfigure}
  \caption{Q4-3 三种策略的费用结构与日费用分布}
  \label{fig:q4_cost_daily}
\end{figure}

由图~\ref{fig:q4_cost} 可见，三种策略的费用差异主要来自紧急购电。N0 的紧急购电费用较高，M0L 通过选择性调整有所降低，而 LA 进一步压低了紧急购电费用，因此全年总费用最低。图~\ref{fig:q4_daily} 进一步表明，LA 的日总费用整体较低，同时高费用日期明显减少。这说明日内滚动调整不仅能够降低全年总费用，还能够减少预测偏差造成的高费用情况。


% 避免三级标题孤悬于页底。
\Needspace{3\baselineskip}
\subsubsection{代表日期结果分析}

为进一步考察两类购电方案在不同季节下的运行结果，选取四个代表日期，其购电量、购电费用及紧急购电量见表~\ref{tab:q4_dates}。

\begin{table}[!htbp]
  \centering
  \caption{波动电价下指定日期的购电结果}
  \label{tab:q4_dates}
  \begingroup
  \small
  \setlength{\tabcolsep}{3pt}
  \begin{tabular}{@{}lrrr@{\hspace{7pt}}rrr@{}}
    \toprule
    日期 & \multicolumn{3}{c}{Q4-2 日前购电方案}
         & \multicolumn{3}{c}{Q4-3 日内滚动方案} \\
    \cmidrule(lr){2-4}\cmidrule(lr){5-7}
    & \shortstack{购电量\\（kWh）}
    & \shortstack{购电费\\（元）}
    & \shortstack{紧急购电量\\（kWh）}
    & \shortstack{购电量\\（kWh）}
    & \shortstack{购电费\\（元）}
    & \shortstack{紧急购电量\\（kWh）} \\
    \midrule
    2025-03-20 & $70\,570.86$ & $44\,640.38$ & 0.00
               & $66\,824.83$ & $44\,186.70$ & 165.68 \\
    2025-06-21 & $35\,205.93$ & $17\,564.67$ & 0.00
               & $33\,936.26$ & $16\,638.47$ & 0.00 \\
    2025-09-23 & $68\,848.07$ & $44\,716.63$ & 272.06
               & $67\,439.16$ & $47\,671.18$ & 631.22 \\
    2025-12-21 & $98\,970.30$ & $75\,354.93$ & 0.00
               & $96\,674.19$ & $75\,675.02$ & 291.73 \\
    \bottomrule
  \end{tabular}
  \endgroup
\end{table}

由表~\ref{tab:q4_dates} 可见，两类方案的购电需求均具有明显的季节差异。以 $6$ 月 $21$ 日和 $12$ 月 $21$ 日为例，Q4-2 的全天购电量分别为 $35\,205.93$ kWh 和 $98\,970.30$ kWh，Q4-3 分别为 $33\,936.26$ kWh 和 $96\,674.19$ kWh，冬季代表日的购电需求明显高于夏季代表日，反映了负载与光伏发电水平的季节变化。

同时，部分代表日期中 Q4-3 的紧急购电量高于 Q4-2。例如 $9$ 月 $23$ 日，两者分别为 $631.22$ kWh 和 $272.06$ kWh，说明日内滚动调整并不能保证每个单日的紧急购电量都低于日前方案。其效果应从全年范围评价：结合表~\ref{tab:q4_all} 可知，LA 相较不进行日内调整的 N0 将全年紧急购电量降低约 $85.6\%$，说明日内调整的主要优势在于整体降低预测偏差造成的紧急购电风险，而非保证所有单日结果均达到最低。

\subsection{模型检验}

为验证问题四调度结果的可靠性，从信息使用和运行约束两方面进行检验。对全年 $334$ 个决策日逐日核验后，未发现未来电价信息泄漏：每日 $0{:}00$ 的电价预测仅使用此前已获得的历史数据，$6{:}00$、$12{:}00$ 和 $18{:}00$ 的日内调整也未使用当天尚未获得的真实电价。

同时，将 Q4-2 与 Q4-3 的全年调度结果重新代入模型进行检查，所有决策日均满足储电量范围 $1200\le E_t\le10800$，且储能接口功率、电量平衡及不可向外网售电等约束均未出现违反情况。按照附件~4 的真实电价重新计算各项费用后，结果与模型输出一致。综上，所得方案符合题目规定的信息条件与运行约束，费用计算结果具有良好的数值一致性。

% 新增问题时复制一个 qx.tex 并在此处加入一行。

\section{综合结果分析}

表~\ref{tab:overall} 汇总各问的正式结果及其对照。问题一为典型日，问题二至四为 $2025$ 年 $2$--$12$ 月共 $334$ 天；费用降幅均相对于本问同口径基准计算。

\begin{table}[!htbp]
  \centering
  \caption{四问关键结果对照}
  \label{tab:overall}
  \small
  \setlength{\tabcolsep}{4pt}
  \begin{tabularx}{\textwidth}{@{}l>{\raggedright\arraybackslash}X>{\raggedright\arraybackslash}X>{\raggedright\arraybackslash}X>{\raggedright\arraybackslash}X@{}}
    \toprule
    问题 & 适用条件与新增机制 & 本文方案 & 关键费用（元） & 同口径对照与节省 \\
    \midrule
    一 & 典型日、确定性日前计划 & 线性规划 & 日购电费 $35\,126.85$ & 无储能 $48\,052.05$，省 $26.90\%$ \\
    \addlinespace
    二 & 全年日前计划；负载/光伏预测；$5$ 倍紧急购电 & 净分位 $+$ 日末 $2400$ $+$ $\rho$ 回放 & 全年 $13\,695\,403.61$ & 固定 $0.8$ 为 $13\,765\,167.58$，省 $0.51\%$ \\
    \addlinespace
    三 & 日内滚动调整；少买 $0.5$ 倍、多买 $1.5$ 倍偏差结算 & $48$ h 展望 $+$ 分源残差修正 & 全年 $13\,275\,795.85$ & 不调整 N0 $15\,197\,639.94$，省 $12.6\%$ \\
    \addlinespace
    四 & 实时波动电价；决策用 $\hat p_0$、结算用附件 4 & 三层电价预测 $+$ 净分位 $+\rho$ / LA $q{=}0.60$ & Q4-2 $14\,413\,479.02$；Q4-3 $14\,001\,199.26$ & 固定 $14\,480\,739.96$，省 $0.46\%$；N0 $15\,983\,608.51$，省 $12.4\%$ \\
    \bottomrule
  \end{tabularx}
\end{table}

由表~\ref{tab:overall} 可见，线性规划在确定性信息下统筹全天调度；预测条件下，日内滚动调整较不调整降低了费用，这一方向在两种电价情景下均成立。各问的信息条件与电价不同，费用绝对值不直接用于比较模型优劣。

参数标定和可行性检验见各问。功率口径对照、旧方案消融及完美信息诊断集中保存在支撑材料 \texttt{condensed\_sections\_notes.md}，其中历史费用不替代本表的正式结果。

\section{模型评价}

\subsection{模型优点}

\textbf{(1) 结构统一、约束对齐题目。} 四个问题都按“预测 $\rightarrow$ 优化 $\rightarrow$ 因果回放”推进：问题一是确定性线性规划，问题二到四是滚动线性规划加锁死合同的因果回放。目标函数统一为购电费，约束与题目逐条对应——储电量运行窗 $1200$--$10800$ kWh、充放电效率 $0.9$、接口功率 $5000$ kW、紧急购电 $5$ 倍、调整部分少买 $0.5$ 倍/多买 $1.5$ 倍。任一时刻的购电、充放电、弃光和紧急购电都能追溯到具体约束。

\textbf{(2) 可复现。} 算法与超参数全部固定（XGBoost 随机种子 $42$、线性规划均求到最优、参数冻结），并配有泄漏检查（预测只用 $D$ 日前信息，错误 $0$ 条）、价格泄漏检查（$0$ 条）与逐时段约束校验（能量平衡、SOC、功率、同时充放电均为 $0$ 错误）。同一套代码重复运行得到同样数字。

\textbf{(3) 有对照、参数在样本外前标定。} 每个问题都设了可对照基线：问题一对照无储能与分时贪婪；问题二对照无分位、固定分量分位，再在净分位上叠加日末约束与保护回放；问题三对照不调整的 N0 与选择门 M0L；问题四对照固定裕度与完美电价 oracle。问题二的 $\alpha$、$\rho$、$S^{*}$ 在 $1$ 月 $15$--$31$ 日网格选定，问题三的 $q_L$ 与展望长度在标定区间选定后再做全年回测。

\subsection{模型不足与推广}

\textbf{(1) 预测依赖历史平稳性。} 负载、光伏与电价都用“历史同刻均值 + XGBoost”类统计模型，特征只有时间，不含天气、节假日与极端事件。样本外出现未见过形态时，偏差会被 $5$ 倍紧急电价放大；结论基于附件 2/附件 4 的 2025 年样本，换年份需重新标定。

\textbf{(2) 若干简化假设。} 充放电效率恒为 $0.9$、储能无自放电与寿命损耗、电价与负载相互独立、$10$ min 内功率恒定；寿命成本未计入目标，实际费用可能因此偏高或偏低。

\textbf{(3) 回放是因果近似而非全知最优。} 日内回放用贪心与滚动 LP，是在“不利用未来信息”下的近似；问题四完美电价诊断显示，距用真实电价决策还差约 $0.57\%$，实际执行偏差也会带来额外费用。

\textbf{(4) 参数在标定区间内选定。} 问题二的 $S^{*}$ 与 $\rho$ 只在同年 $1$ 月 $16$ 格选定，问题三的分位与缓冲系数在各自区间内选定，并非独立年份的样本外检验。可改进方向：给预测加天气特征、用随机规划或鲁棒优化刻画尾部风险、把储能损耗计入目标，并做多年份滚动验证。

\section*{AI工具使用声明}
本参赛队在竞赛过程中使用了 AI 工具，主要用于\textbf{在此填写实际用途，例如语言润色、代码调试或格式检查}，详细使用情况见支撑材料。\par
% 若未使用 AI，请将上句替换为：
% 本参赛队在竞赛过程中未使用任何AI工具。

\begin{thebibliography}{99}

\bibitem{lasseter2002}
Lasseter R H. MicroGrids[C]//2002 IEEE Power Engineering Society Winter Meeting. IEEE, 2002, 1: 305-308.

\bibitem{hatziargyriou2007}
Hatziargyriou N, Asano H, Iravani R, et al. Microgrids[J]. IEEE Power and Energy Magazine, 2007, 5(4): 78-94.

\bibitem{katiraei2008}
Katiraei F, Iravani R, Hatziargyriou N, et al. Microgrids management[J]. IEEE Power and Energy Magazine, 2008, 6(3): 54-65.

\bibitem{dunn2011}
Dunn B, Kamath H, Tarascon J M. Electrical energy storage for the grid: A battery of choices[J]. Science, 2011, 334(6058): 928-935.

\bibitem{chen2009}
Chen H, Cong T N, Yang W, et al. Progress in electrical energy storage system: A critical review[J]. Progress in Natural Science, 2009, 19(3): 291-312.

\bibitem{boyd2004}
Boyd S, Vandenberghe L. Convex Optimization[M]. Cambridge: Cambridge University Press, 2004.

\bibitem{conejo2010}
Conejo A J, Carrión M, Morales J M. Decision Making Under Uncertainty in Electricity Markets[M]. New York: Springer, 2010.

\bibitem{chen2016xgboost}
Chen T, Guestrin C. XGBoost: A Scalable Tree Boosting System[C]//Proceedings of the 22nd ACM SIGKDD International Conference on Knowledge Discovery and Data Mining. ACM, 2016: 785-794.

\bibitem{taylor2003}
Taylor J W. Short-term electricity demand forecasting using double seasonal exponential smoothing[J]. Journal of the Operational Research Society, 2003, 54(8): 799-805.

\bibitem{conejo2005}
Conejo A J, Contreras J, Espínola R, et al. Forecasting electricity prices for a day-ahead pool-based electric energy market[J]. International Journal of Forecasting, 2005, 21(3): 435-462.

\bibitem{parisio2014}
Parisio A, Rikos E, Glielmo L. A Model Predictive Control Approach to Microgrid Operation Optimization[J]. IEEE Transactions on Control Systems Technology, 2014, 22(5): 1813-1827.

\bibitem{farzin2017}
Farzin H, Fotuhi-Firuzabad M, Moeini-Aghtaie M. A Stochastic Multi-Objective Framework for Optimal Scheduling of Energy Storage Systems in Microgrids[J]. IEEE Transactions on Smart Grid, 2017, 8(1): 117-127.

\end{thebibliography}

\begin{appendices}
\section{支撑材料文件列表}
列出提交的支撑材料文件名及其用途。附录中不得出现队员、学校或赛区身份信息。

\section{关键代码}
完整、可运行的源程序应作为支撑材料提交；篇幅较短的关键代码可以放在此处。\par
% \lstinputlisting[language=Python]{../code/Q1/main.py}
\end{appendices}
